\subsection{INVERSE LIMITS OF UNIFORM SPACES}

Let I be a partially ordered set in which the partial ordering is written \(\alpha \leqslant \beta\) . For each \(\alpha \in I\) let \(X_{\alpha}\) be a uniform space, and for each pair of indices \(\alpha, \beta\) such that \(\alpha \leqslant \beta\) let \(f_{\alpha \beta}\) be a mapping of \(X_{\beta}\) into \(X_{\alpha}\) .

We shall say that \((\mathbf{X}_{\alpha}, f_{\alpha\beta})\) is an inverse system of uniform spaces if (i) \((\mathbf{X}_{\alpha}, f_{\alpha\beta})\) is an inverse system of sets (cf.~Chapter I, Appendix, no. 1) and (ii) whenever \(\alpha \leqslant \beta\) , \(f_{\alpha\beta}\) is uniformly continuous. On \(X = \varprojlim X_{\alpha}\) the coarsest uniformity for which the canonical mappings \(f_{\alpha}: X \to X_{\alpha}\) are uniformly continuous is called the inverse limit (with respect to the \(f_{\alpha\beta}\) ) of the uniformities of the \(X_{\alpha}\) , and the set X endowed with this coarsest uniformity is called the inverse limit of the inverse system of uniform spaces \((\mathbf{X}_{\alpha}, f_{\alpha\beta})\) . All the properties of inverse limits of topological spaces established in Chapter I, § 4, no. 4 (with the exception of Proposition 9) remain valid if we replace ``topology'' by ``uniformity'' and ``continuous mapping'' by ``uniformly continuous mapping''. In addition:

PROPOSITION 10. Let I be a directed set, let \((\mathbf{X}_{\alpha}, f_{\alpha \beta})\) be an inverse system of uniform spaces indexed by I, and let J be a cofinal subset of I. For each \(\alpha \in I\) let \(f_{\alpha}\) be the canonical mapping of \(\mathbf{X} = \varprojlim \mathbf{X}_{\alpha}\) into \(\mathbf{X}_{\alpha}\) , and let \(g_{\alpha}\) denote \(f_{\alpha} \times f_{\alpha}\) . Then the family of sets \(\overline{g}_{\alpha}^{1}(\mathrm{V}_{\alpha})\) , where \(\alpha\) runs through J and where, for each \(\alpha \in J\) , \(\mathrm{V}_{\alpha}\) runs through a fundamental system of entourages of \(\mathrm{X}_{\alpha}\) , is a fundamental system of entourages of X.

We leave the proof to the reader; it is a straightforward adaptation of the proof of Proposition 9 of Chapter I, § 4, no. 4.

Finally, the topology on \(X = \varprojlim X_{\alpha}\) induced by the inverse limit of the uniformities of the \(X_{\alpha}\) is the inverse limit of the topologies of the \(X_{\alpha}\) .

